\documentclass[instructions.tex]{subfiles}
\begin{document}

\chapter{On risque son salut, quand on diffère de quitter le péché et de se convertir.}


Pour deux raisons : parce qu'on peut être surpris de la mort, et qu'en différant on est en effet ordinairement surpris.

\primo On peut être surpris, personne n'en doute ; on peut l'être à tout moment. Il vaudrait donc mieux souffrir pendant cinquante ans les feux et les maladies les plus aiguës, que de rester seulement un quart d'heure en péché mortel. Cinquante ans de tourmens au bout de quelques années seraient finis ; mais pendant le quart d'heure que vous restez dans le péché mortel, si Dieu prononçait votre sentence, votre malheur serait sans fin.

Deux momens peuvent décider de votre éternité. Dans le moment présent vous pouvez vous convertir à Dieu, et vous mettre en état d'être heureux pour jamais ; mais le moment suivant, vous ne le pourrez peut-être plus ; et si la mort vous surprend dans ce moment, votre perte sera sans ressource. Si vous différez de rompre les chaînes qui vous attachent au péché, il ne faudra qu'un moment pour les changer en des chaînes de feu que vous ne pourrez plus briser, puisqu'elles seront éternelles. \emph{Ne tardez donc pas}, dit le Sage, \emph{de retourner à Dieu pendant qu'il est temps, car il vous surprendra tout-à-coup dans sa colère}.

\secundo Qui ne tremblerait pas à cette menace : \emph{Dieu vous surprendra dans sa colère} ? car si Dieu est fidèle dans ses promesses, il ne l'est pas moins dans ses menaces. S'il attend le pécheur avec patience, il le punit aussi avec force\footnote{Judex patiens et fortis.}. Il punit de la sorte quand on y pense le moins. Un homme qui veut se venger d'un ennemi ne saisit pas toujours la première occasion ; il attend une conjoncture plus propre à ses desseins ; jusqu'à ce qu'ayant trouvé le moment, il décharge sur lui sa colère.

Dieu pourrait punir l'homme aussitôt après son péché ; mais il attend souvent l'occasion de le punir avec plus d'éclat. Pendant ce temps Dieu semble le combler de prospérités ; ce pécheur jouit des honneurs, des biens de la vie : mais Dieu voit le moment auquel il va surprendre ce misérable, et le faire passer, par une mort imprévue, au dernier des malheurs.

Ainsi punit-il autrefois l'impudente Jézabel, l'impie Antiochus ; ainsi punit-il le roi Balthasar. Ce roi, dans un festin, fit apporter les vases sacrés pour y boire à l'honneur de ses faux dieux. C'était à ce sacrilège que Dieu l'attendait pour se venger. Au milieu du festin on vit une main qui écrivait sur la muraille l'arrêt de sa mort. À ce spectacle, le roi fut rempli d'une frayeur mortelle en présence de ses courtisans. Dieu l'avait ainsi disposé, afin que ceux qui étaient les complices de ses crimes, fussent aussi les témoins de sa punition, châtiment qui suivit de près : il fut massacré la même nuit.

Dieu ne vous punira peut-être pas d'une manière éclatante, comme ce malheureux prince ; mais il vous punira d'une manière aussi terrible, en vous surprenant par une mort funeste au milieu de vos désordres. Hommes aveugles, vous ne voudriez pas mourir dans le crime ! pourquoi donc y persévérez-vous ? \emph{C'est le propre du démon}, dit saint Bernard, \emph{de persévérer dans le mal ; ceux qui y persévèrent méritent de périr avec lui}. Revenez donc promptement à Dieu, de crainte que sa main ne s'appesantisse sur vous, et ne vous ôte le temps de le faire. Le mal devient incurable quand on néglige le remède.


\section{Résolutions}

\primo Enfin je ne veux plus risquer mon salut.

\secundo En conséquence, il faut que ce jour même soit l'époque fixe d'une nouvelle vie pour moi. 

Mon Dieu, convertissez-moi, et je serai converti ; je compte plus sur votre grâce que sur mes résolutions, dont je reconnais, à ma honte, l'inconstance et la faiblesse extrême.

\end{document}
